% Provenance: C_F0031 and C_F0068 Part C, finite-rank closure and domains.
\section{Dual system criterion}
\label{sec:dual-system-criterion}

We prove Theorem~\ref{thm:analytic} by applying the liberation dual system criterion of Theorem~\ref{thm:intrinsic-analytic}.
We first prove that the existence of a finite-rank closure and its rank are independent of the fraction pair, and identify finite-rank solutions of equation~\eqref{eq:certificate} with the closures of the corresponding fraction commutators.
We then construct a normalized liberation dual system directly from the dual operators and apply assertion~\textup{(\ref{item:intrinsic-analytic-fraction-commutators})} of Theorem~\ref{thm:intrinsic-analytic}.

The following lemma supplies the domain facts needed for the fraction commutators.
\begin{lemma}\label{lem:fraction-domains}
Let \(D_1,\ldots,D_d\) be densely defined linear operators on \(H\) such that
\[\M\Omega\subseteq\dom D_i\cap\dom D_i^*\]
for \(1\leq i\leq d\).
For every \(u\in L^0(\M)\), every fraction pair \(Q\) of \(u\), and every \(i\), the operator \(\Gamma[Q,D_i]\) is closable and densely defined.
\end{lemma}
\begin{proof}
Fix \(i\) and write \(Q=(s,a,b,t)\).
Since \(b,t\in\M\) preserve \(\M\Omega\), the domain of \(\Gamma[Q,D_i]\) contains \(\M\Omega\).
For \(\eta\in\M\Omega\) and \(\xi\in\dom\Gamma[Q,D_i]\), the hypothesis of the lemma gives
\begin{equation}\label{eq:defect-adjoint}
 \langle\eta,\Gamma[Q,D_i]\xi\rangle
 =\left\langle
 b^*D_i^*(s^*\eta)-t^*D_i^*(a^*\eta),\xi
 \right\rangle.
\end{equation}
Thus \(\M\Omega\subseteq\dom(\Gamma[Q,D_i])^*\), so \(\Gamma[Q,D_i]\) is closable.
\end{proof}

The next lemma allows us to pass from one fraction pair to every fraction pair before applying Theorem~\ref{thm:intrinsic-analytic}.

\begin{lemma}\label{lem:fraction-ranks}
Let $D$ be a densely defined linear operator on $H$ such that
$\M\Omega\subseteq\dom D\cap\dom D^*$.
For $u\in L^0(\M)$ and any two fraction pairs $Q,Q'$ for $u$,
\[
 \overline{\Gamma[Q,D]}\in\F(H)
 \quad\Longleftrightarrow\quad
 \overline{\Gamma[Q',D]}\in\F(H).
\]
When these conditions hold, the two closures have the same rank.
\end{lemma}
\begin{proof}
Write $Q=(s,a,b,t)$ and $Q'=(s',a',b',t')$.
Lemma~\ref{lem:common} gives regular $x,y,p,q\in\M$ such that $xs=ys'$ and $tp=t'q$.
Since $xa=ya'$ and $bp=b'q$, we have on $\M\Omega$
\begin{equation}\label{eq:fraction-common-comparison}
 x\Gamma[Q,D]p=y\Gamma[Q',D]q.
\end{equation}
Suppose that $\overline{\Gamma[Q,D]}\in\F(H)$.
The left-hand side of equation~\eqref{eq:fraction-common-comparison} has finite rank, and $y$ is injective.
Thus the restriction of $\Gamma[Q',D]$ to the dense subspace $q\M\Omega$ has finite rank.
It is closable by Lemma~\ref{lem:fraction-domains}, so Lemma~\ref{lem:finite-closure} makes its closure a bounded finite-rank operator on all of $H$.
The closure of $\Gamma[Q',D]$ extends this operator and hence equals it.
Interchanging $Q$ and $Q'$ proves the equivalence.

Both closures are now bounded, so equation~\eqref{eq:fraction-common-comparison} extends to
\[
 x\overline{\Gamma[Q,D]}p=y\overline{\Gamma[Q',D]}q
 \qquad\text{on }H.
\]
The left factors $x,y$ are injective, while $p,q$ have dense range.
These multiplications preserve the rank of a bounded finite-rank operator, proving the rank equality.
\end{proof}

We next identify any finite-rank solution of equation~\eqref{eq:certificate} with the fraction commutators of a dual system.

\begin{lemma}\label{lem:finite-rank-identification}
Let $X=(X_1,\ldots,X_d)\in\M^d$, and let $(D_1,\ldots,D_d)$ be a dual system for $X$.
Let $Q$ be a fraction pair for $u\in L^0(\M)$.
If $T_1,\ldots,T_d\in\F(H)$ satisfy
\[
 \sum_j[\rho(X_j),T_j]=\Gamma[Q,P_\Omega],
\]
then $\overline{\Gamma[Q,D_i]}=T_i$ for every $i$.
\end{lemma}
\begin{proof}
Write $Q=(s,a,b,t)$ and put $\ell_j:=-\Theta^{-1}(T_j)\in H\otimes_{\C}H$.
By Lemma~\ref{lem:tensor-certificate},
\[
 \sum_j\Psi_j(\ell_j)=a\otimes t^{\mathrm{op}}-s\otimes b^{\mathrm{op}}.
\]
Lemma~\ref{lem:tensor-clearing} gives regular $s',t'\in\M$ such that
$\ell_j':=(s'\otimes(t')^{\mathrm{op}})\ell_j\in\M\otimes_{\C}\M^{\mathrm{op}}$ for every $j$.
Equation~\eqref{eq:psi-left-linearity} gives
\[
 \sum_j\Psi_j(\ell_j')
 =s'a\otimes(tt')^{\mathrm{op}}-s's\otimes(bt')^{\mathrm{op}}.
\]
For fixed $i$ and $\xi\in\M\Omega$, apply to this equality the linear map
\[
 \M\otimes_{\C}\M^{\mathrm{op}}\longrightarrow H,\qquad
 c\otimes e^{\mathrm{op}}\longmapsto cD_i(e\xi).
\]
All inputs of $D_i$ lie in $\M\Omega$, and the dual relations give
\[
 cX_jD_i(e\xi)-cD_i(X_je\xi)=\delta_{ij}cP_\Omega e\xi.
\]
Using equation~\eqref{eq:theta-bounded}, we obtain
\[
 -s'\Gamma[Q,D_i]t'\xi
 =\Theta(\ell_i')\xi=-s'T_it'\xi.
\]
Since $s'$ is injective, $\Gamma[Q,D_i]$ agrees with $T_i$ on the dense subspace $t'\M\Omega$.
Lemma~\ref{lem:fraction-domains} gives closability of $\Gamma[Q,D_i]$.
Its closure therefore extends the bounded operator $T_i$, whose domain is all of $H$, and hence equals $T_i$.
\end{proof}

For \(N,S\in\B(H)\) with \(N\) normal, the classical Fuglede theorem \cite[Theorem~I, p.~35]{Fug50} gives \(SN=NS\implies SN^*=N^*S\).
The following analogue for bounded linear maps from \(\M\) to \(H\) will be used in Lemma~\ref{lem:module-commutation}.

\begin{lemma}[{see \cite[Proposition~1]{CS77}}]\label{lem:normal-intertwining}
Let \(x\in\M\) be normal, and let \(T:\M\to H\) be a bounded linear map.
If \(T(xz)=xT(z)\) for every \(z\in\M\), then \(T(x^*z)=x^*T(z)\) for every \(z\in\M\).
\end{lemma}

\begin{lemma}\label{lem:module-commutation}
Let \(T:\M\Omega\to H\) be a closable linear operator.
Then \(z\mapsto T(z\Omega)\) is bounded from \(\M\) to \(H\).
If \(x_1,\ldots,x_m\in\M\) are normal and \(T(x_j\xi)=x_jT\xi\) for every \(j\) and \(\xi\in\M\Omega\), then
\[
 T(a\xi)=aT\xi
 \qquad(a\in W^*(x_1,\ldots,x_m),\ \xi\in\M\Omega).
\]
\end{lemma}
\begin{proof}
Identify \(\M\Omega\) with \(\M\).
The graph of \(T:\M\to H\) is norm closed because the Hilbert-space closure \(\overline T\) is closed and agrees with \(T\) on \(\M\Omega\).
The closed graph theorem therefore gives the asserted bound.
Lemma~\ref{lem:normal-intertwining} gives \(T(x_j^*\xi)=x_j^*T\xi\).
Products and linear combinations give the same relation for the unital $*$-algebra generated by \(x_1,\ldots,x_m\).
For \(a\in W^*(x_1,\ldots,x_m)\), Kaplansky density \cite[Theorem~2.3.1]{AP18} gives a uniformly bounded net \(a_\lambda\) in this algebra converging strongly to \(a\).
For fixed \(\xi\in\M\Omega\),
\[
 a_\lambda\xi\longrightarrow a\xi,\qquad
 T(a_\lambda\xi)=a_\lambda T\xi\longrightarrow aT\xi.
\]
Closedness of \(\overline T\) and \(a\xi\in\M\Omega\) give \(T(a\xi)=aT\xi\).
\end{proof}

% Provenance: Yin_C_F0024-P1, Sections 2--6, and Yin_C_F0033-P1
% specialized to M Omega; the proof of the criterion uses C_F0024-P3
% (v2), identifying the C_F0063 certificates with actual fraction commutators.

We now construct a normalized liberation dual system from a dual system for normal generators and compare their fraction commutators.

\begin{lemma}\label{lem:fundamental}
Let $X_1,\ldots,X_d\in\M$ be normal operators generating $\M$, and let $(D_1,\ldots,D_d)$ be a dual system for $X$.
Then every $X_i-\lambda$, $\lambda\in\C$, is regular.
Define on $\M\Omega$
\begin{equation}\label{eq:dual-liberation-construction}
 K_i:=[\rho(X_i),D_i|_{\M\Omega}],\qquad
 C:=\sum_jK_j-P_\Omega,\qquad
 L_i:=K_i-\frac1d C.
\end{equation}
The tuple $(L_1,\ldots,L_d)$ is a normalized liberation dual system for $X$.
For every fraction pair $Q$ and every $i$, if $\overline{\Gamma[Q,D_i]}$ is bounded on $H$, then
\begin{equation}\label{eq:dual-liberation-fractions}
 \overline{\Gamma[Q,L_i]}
 =[\rho(X_i),\overline{\Gamma[Q,D_i]}].
\end{equation}
In particular, if $\overline{\Gamma[Q,D_i]}\in\F(H)$ for every $i$, then
\begin{equation}\label{eq:fundamental}
 \sum_{i=1}^d[\rho(X_i),\overline{\Gamma[Q,D_i]}]=\Gamma[Q,P_\Omega].
\end{equation}
\end{lemma}
\begin{proof}
For a spectral projection $e:=1_{\{\lambda\}}(X_i)$, the dual relation gives
$eP_\Omega e=e[X_i,D_i]e=0$ on $\M\Omega$.
Density and faithfulness of the trace give $e=0$.
Normality then implies that $X_i-\lambda$ is regular.

Since $\rho(X_i)^*=\rho(X_i^*)$ preserves $\M\Omega$, the adjoint-domain hypothesis gives
\[
 \M\Omega\subseteq\dom K_i^*,\qquad
 K_i^*\eta=D_i^*\rho(X_i^*)\eta-\rho(X_i^*)D_i^*\eta
 \quad(\eta\in\M\Omega).
\]
Consequently, $\M\Omega$ is contained in the adjoint domains of $C$ and each $L_i$ as well.
Left and right multipliers commute, and traciality gives
$[\rho(X_i),P_\Omega]=[X_i,P_\Omega]$ on $H$.
Jacobi's identity therefore gives on $\M\Omega$
\[
 [X_j,K_i]=[\rho(X_i),[X_j,D_i]]
 =\delta_{ij}[X_i,P_\Omega],\qquad [X_j,C]=0.
\]
It follows that $[X_j,L_i]=\delta_{ij}[X_i,P_\Omega]$ and $\sum_iL_i=P_\Omega$ on $\M\Omega$.
Thus $(L_1,\ldots,L_d)$ is a normalized liberation dual system.

The operator $C$ is closable, so Lemma~\ref{lem:module-commutation} gives
$C(a\xi)=aC\xi$ for all $a\in\M$ and $\xi\in\M\Omega$.
For $Q=(s,a,b,t)$, this implies
\[
 \Gamma[Q,C]\xi=(sb-at)C\xi=0\qquad(\xi\in\M\Omega).
\]
Right multipliers commute with the left multipliers in $Q$, and hence
\begin{equation}\label{eq:dual-liberation-on-domain}
 \Gamma[Q,L_i]\xi
 =\Gamma[Q,K_i]\xi
 =\rho(X_i)\Gamma[Q,D_i]\xi-\Gamma[Q,D_i]\rho(X_i)\xi
 \qquad(\xi\in\M\Omega).
\end{equation}
If $\overline{\Gamma[Q,D_i]}$ is bounded on $H$, the right-hand side extends to
$[\rho(X_i),\overline{\Gamma[Q,D_i]}]\in\B(H)$.
By Lemma~\ref{lem:fraction-domains}, $\Gamma[Q,L_i]$ is closable on its natural domain.
Its closure extends this bounded operator on all of $H$, proving equation~\eqref{eq:dual-liberation-fractions}.
Finally, summing equation~\eqref{eq:dual-liberation-fractions} on $\M\Omega$ and using $\sum_iL_i=P_\Omega$ gives equation~\eqref{eq:fundamental} by density.
\end{proof}

\begin{proof}[Proof of Theorem~\ref{thm:analytic}]
Let $(L_1,\ldots,L_d)$ be the normalized liberation dual system constructed in Lemma~\ref{lem:fundamental}.
That lemma also gives regularity of every scalar shift, so Theorem~\ref{thm:intrinsic-analytic} applies.
Fix $u\in L^0(\M)$ and a fraction pair $Q$ for $u$.

$\textup{(1)}\implies\textup{(3)}$.
By Theorem~\ref{thm:intrinsic-analytic}\textup{(\ref{item:intrinsic-analytic-fraction-commutators})}, there are $T_1,\ldots,T_d\in\F(H)$ such that
\[
 \overline{\Gamma[Q,L_i]}=[\rho(X_i),T_i]\qquad(1\leq i\leq d).
\]
Summing on $\M\Omega$ and using $\sum_iL_i=P_\Omega$, then extending by boundedness, gives
\[
 \sum_i[\rho(X_i),T_i]=\Gamma[Q,P_\Omega].
\]
Lemma~\ref{lem:finite-rank-identification} gives $\overline{\Gamma[Q,D_i]}=T_i$ for every $i$, proving assertion~\textup{(3)}.

$\textup{(3)}\implies\textup{(1)}$.
By Lemma~\ref{lem:fraction-ranks}, $\overline{\Gamma[Q',D_i]}\in\F(H)$ for every fraction pair $Q'$ of $u$ and every $i$.
Equation~\eqref{eq:dual-liberation-fractions} therefore gives
\[
 \overline{\Gamma[Q',L_i]}
 =[\rho(X_i),\overline{\Gamma[Q',D_i]}]
 \in[\rho(X_i),\F(H)].
\]
Thus assertion~\textup{(\ref{item:intrinsic-analytic-fraction-commutators})} of Theorem~\ref{thm:intrinsic-analytic} holds, and $u\in\Div_{L^0(\M)}\A_X$.

$\textup{(2)}\Longleftrightarrow\textup{(3)}$.
Lemma~\ref{lem:fraction-domains} gives closability, and Lemma~\ref{lem:finite-closure} gives the forward implication.
The reverse implication follows by restriction.
The rank independence is Lemma~\ref{lem:fraction-ranks}.
\end{proof}
